% !TEX program = xelatex
% Version grands caractères (paysage, 22pt) du cours de la Séquence 2
% (Vecteurs), pour une élève déficiente visuelle. Document autonome
% (n'utilise PAS CLT2627.sty), sur le même principe que "Séquence 1 -
% Cours - Calculs Numériques - Grands Caractères.tex" : une seule colonne
% pleine largeur, texte non justifié, interligne large, gras plutôt
% qu'italique pour les emphases.
%
% Contrairement à la version Séquence 1 : ce chapitre étant très
% géométrique, TOUT le contenu du cours original est conservé, y compris
% les schémas (TikZ, redessinés à plus grande échelle pour rester lisibles
% en 22pt) et les cartes vidéos Yvan Monka (icône + titre cliquable + QR
% code, agrandis). Seuls les tableaux "Objectif / exercices" (renvois vers
% le manuel et la feuille d'exercices, hors sujet pour une version cours)
% sont retirés, comme dans la version Séquence 1.
%
% VERSION COMPLÉTÉE : copie exacte de "Séquence 2 - Cours - Vecteurs -
% Grands Caractères.tex" (même découpage page par page, un bloc par page),
% où l'on révèle ce que l'élève doit compléter, en reprenant les réponses
% de "Séquence 2 - Cours - Vecteurs - Version complétée.tex" (mise en
% page classique). Convention de couleurs, la même que dans cette version
% classique :
%   - les mots clés déjà donnés dans la version élève (en rouge
%     \surligneReponse, ex-J1 du poly) restent tels quels ;
%   - tout ce que l'élève complète (réponses des Exemples cachées côté
%     élève par \Trou / \phantom / tracés TikZ invisibles, cases à
%     cocher, constructions sur quadrillage) est en vert un peu foncé
%     (\vertComplet), pour le distinguer d'un coup d'œil du cours. Les
%     vecteurs intermédiaires d'une construction bout à bout sont en
%     tirets verts, le vecteur somme en rouge plein (étape 4 de la
%     méthode).
% Pour reprendre ce fichier après une modification côté élève : repartir
% du fichier élève et reporter les blocs \Complet / vertComplet.
\documentclass[a4paper,landscape]{article}

\usepackage{fontspec}
\usepackage[french]{babel}
\usepackage{geometry}
\geometry{a4paper,landscape,top=1.4cm,bottom=2cm,left=1.8cm,right=1.8cm,footskip=1.2cm}

\setmainfont{Helvetica Neue} % police sans-serif, lisible en grands caractères

\usepackage{amsmath,amssymb}
\usepackage{xcolor}
\usepackage[most]{tcolorbox}
\usepackage{enumitem}
\usepackage{ragged2e}
\usepackage{fancyhdr}
\usepackage[hidelinks]{hyperref}
\usepackage{lastpage}
\usepackage{tikz}
\usepackage{multicol}
\usepackage{fontawesome}
\usepackage{qrcode}
\usepackage{titlesec}
\usepackage{frcursive}
\usetikzlibrary{arrows.meta,shapes.misc,shapes.geometric}

% Croix de placement des points (identique au cours version normale, cf.
% "Séquence 2 - Cours - Vecteurs.tex"), simplement agrandie pour le format
% grands caractères.
\tikzset{ptcross/.style={cross out,draw,very thick,minimum size=3mm,inner sep=0pt}}

% ============================================================ %
% Correctif du bug de crénage de frcursive sous XeLaTeX (repris de
% CLT2627.sty) : sans lui, certaines paires de lettres consécutives dans
% \cursive explosent (une lettre projetée loin à droite, hors de la
% page), notamment "lettre + s".
% ============================================================ %
\makeatletter
\ifdefined\XeTeXcharclass
\XeTeXcharclass `a 201 \XeTeXcharclass `b 201 \XeTeXcharclass `c 201
\XeTeXcharclass `d 201 \XeTeXcharclass `e 201 \XeTeXcharclass `f 201
\XeTeXcharclass `g 201 \XeTeXcharclass `h 201 \XeTeXcharclass `i 201
\XeTeXcharclass `j 201 \XeTeXcharclass `k 201 \XeTeXcharclass `l 201
\XeTeXcharclass `m 201 \XeTeXcharclass `n 201 \XeTeXcharclass `o 201
\XeTeXcharclass `p 201 \XeTeXcharclass `q 201 \XeTeXcharclass `r 201
\XeTeXcharclass `s 201 \XeTeXcharclass `t 201 \XeTeXcharclass `u 201
\XeTeXcharclass `v 201 \XeTeXcharclass `w 201 \XeTeXcharclass `x 201
\XeTeXcharclass `y 201 \XeTeXcharclass `z 201
\XeTeXcharclass `A 201 \XeTeXcharclass `B 201 \XeTeXcharclass `C 201
\XeTeXcharclass `D 201 \XeTeXcharclass `E 201 \XeTeXcharclass `F 201
\XeTeXcharclass `G 201 \XeTeXcharclass `H 201 \XeTeXcharclass `I 201
\XeTeXcharclass `J 201 \XeTeXcharclass `K 201 \XeTeXcharclass `L 201
\XeTeXcharclass `M 201 \XeTeXcharclass `N 201 \XeTeXcharclass `O 201
\XeTeXcharclass `P 201 \XeTeXcharclass `Q 201 \XeTeXcharclass `R 201
\XeTeXcharclass `S 201 \XeTeXcharclass `T 201 \XeTeXcharclass `U 201
\XeTeXcharclass `V 201 \XeTeXcharclass `W 201 \XeTeXcharclass `X 201
\XeTeXcharclass `Y 201 \XeTeXcharclass `Z 201
\XeTeXcharclass `\é 201 \XeTeXcharclass `\è 201 \XeTeXcharclass `\ê 201
\XeTeXcharclass `\ë 201 \XeTeXcharclass `\É 201 \XeTeXcharclass `\È 201
\XeTeXcharclass `\Ê 201 \XeTeXcharclass `\Ë 201 \XeTeXcharclass `\à 201
\XeTeXcharclass `\â 201 \XeTeXcharclass `\ä 201 \XeTeXcharclass `\À 201
\XeTeXcharclass `\Â 201 \XeTeXcharclass `\Ä 201 \XeTeXcharclass `\ô 201
\XeTeXcharclass `\ö 201 \XeTeXcharclass `\Ô 201 \XeTeXcharclass `\Ö 201
\XeTeXcharclass `\î 201 \XeTeXcharclass `\ï 201 \XeTeXcharclass `\Î 201
\XeTeXcharclass `\Ï 201 \XeTeXcharclass `\ù 201 \XeTeXcharclass `\û 201
\XeTeXcharclass `\ü 201 \XeTeXcharclass `\Ù 201 \XeTeXcharclass `\Û 201
\XeTeXcharclass `\Ü 201 \XeTeXcharclass `\ç 201 \XeTeXcharclass `\Ç 201
\XeTeXcharclass `\œ 201 \XeTeXcharclass `\Œ 201 \XeTeXcharclass `\æ 201
\XeTeXcharclass `\Æ 201
\let\cursive@orig\cursive
\renewcommand{\cursive}{%
  \cursive@orig
  \XeTeXinterchartoks 201 201 = {\kern0pt }%
}
\fi
\makeatother

% ============================================================ %
% Taille de police : 22pt pour le texte courant, tailles voisines
% redéfinies en proportion (utilisées par \section, \tcbtitle, etc.)
% ============================================================ %
\renewcommand{\normalsize}{\fontsize{22}{29}\selectfont}
\renewcommand{\small}{\fontsize{19}{25}\selectfont}
\renewcommand{\large}{\fontsize{25}{31}\selectfont}
\renewcommand{\Large}{\fontsize{29}{35}\selectfont}
\renewcommand{\LARGE}{\fontsize{33}{40}\selectfont}
\renewcommand{\huge}{\fontsize{37}{44}\selectfont}
\renewcommand{\Huge}{\fontsize{41}{48}\selectfont}
\normalsize
\linespread{1.15}
\setlength{\parskip}{10pt}
\setlength{\parindent}{0pt}
\RaggedRight % texte non justifié (ragged2e) dans tout le document

\setlist{leftmargin=1.4em, itemsep=6pt, topsep=4pt}

% Notation vecteur, reprise de CLT2627.sty (\Vj{AB} = vecteur AB).
\newcommand{\Vj}[1]{\overrightarrow{\text{#1}}}

% ============================================================ %
% Titre de séquence (repris tel quel de "Séquence 2 - Cours -
% Vecteurs.tex", macro autonome, indépendante de CLT2627.sty) : coins
% arrondis, étiquette posée via overlay qui entaille le cadre en haut à
% gauche, grand titre centré en dessous. #1 = largeur (fraction de
% \linewidth), #2 = étiquette, #3 = titre.
% ============================================================ %
\newcommand{\SequenceTitre}[3]{%
  \begin{center}
  \begin{tcolorbox}[
    enhanced,
    width=#1\linewidth,
    colback=white,
    colframe=black,
    boxrule=1pt,
    arc=10pt,
    sharp corners=downhill,
    left=14pt, right=14pt, top=12pt, bottom=10pt,
    halign=center,
    overlay={
      \node[fill=white, anchor=west, font=\Large, inner xsep=6pt]
        at ([xshift=16pt]frame.north west) {#2};
    },
  ]
  {\Huge #3}
  \end{tcolorbox}
  \end{center}
}

% ============================================================ %
% Style des sections / sous-sections : repris de CLT2627.sty (même
% présentation que la version classique) — pastille grise numérotée +
% titre en écriture \cursive pour les sections, pastille magenta lettrée
% + titre \cursive magenta pour les sous-sections — simplement agrandi en
% proportion du reste du document (badges et polices redimensionnés pour
% le format 22pt).
% ============================================================ %
\def\thesection{\arabic{section}}
\titlespacing*{\section}{0pt}{4.5ex plus 1ex minus .2ex}{3ex plus .2ex}
\titleformat{\section}[hang]
    {\cursive\Large\color{black}\hyphenchar\font=-1}
    {\begin{tikzpicture}[baseline=(nodetext.base)]
        \node[
            draw=black,
            fill=gray,
            text=white,
            rounded rectangle,
            inner sep=9pt,
            font=\cursive\Large
        ] (nodetext) {\hspace*{0.3em}\thesection\hspace*{0.3em}};
     \end{tikzpicture}}
    {14pt}
    {}

\def\thesubsection{\alph{subsection}}
\titlespacing*{\subsection}{2.6cm}{0.5cm}{0.2cm}
\titleformat{\subsection}[hang]
    {\cursive\large\color{magenta}\hyphenchar\font=-1}
    {\begin{tikzpicture}[baseline=(nodetext.base)]
        \node[
        draw=black,
        fill=magenta!30,
        text=white,
        rounded corners=4pt,
        inner sep=5pt,
        font=\bfseries\normalsize] (nodetext) {\thesubsection};
     \end{tikzpicture}}
    {14pt}
    {}

% ============================================================ %
% Couleurs (reprises de CLT2627.sty pour garder le même code couleur
% Définition/Propriété/Théorème/Exemple/À noter/Attention que le poly)
% ============================================================ %
\definecolor{vertDefinition}{RGB}{46,160,67}
\definecolor{bleuPropriete}{RGB}{41,121,204}
\definecolor{bordeauxTheoreme}{RGB}{196,30,58}
\definecolor{grisExemple}{RGB}{85,85,95}
\definecolor{ambreNoter}{RGB}{214,144,20}
\definecolor{rougeAttention}{RGB}{201,42,42}
\definecolor{surligneReponse}{RGB}{192,42,42} % couleur des réponses (ex-J1)

% ============================================================ %
% Boîtes pleine largeur, "breakable" (peuvent se couper sur une page
% suivante sans risque, puisqu'on n'est plus dans une minipage étroite).
% ============================================================ %
\newcounter{definition}
\newcounter{propriete}
\newcounter{theoreme}
\newcounter{exemple}

\newtcolorbox{BoiteDfn}[1][]{
  enhanced, breakable,
  colback=vertDefinition!6, colframe=vertDefinition,
  boxrule=0pt, borderline west={3pt}{0pt}{vertDefinition},
  arc=0pt, outer arc=0pt,
  left=14pt, right=14pt, top=10pt, bottom=10pt,
  before skip=12pt, after skip=12pt,
}
\newcommand{\Dfn}[2][]{%
  \begin{BoiteDfn}
    \refstepcounter{definition}%
    \textbf{\color{vertDefinition}Définition~\arabic{definition}}%
    \ifx\\#1\\\else\ --~\textbf{#1}\fi
    \par\smallskip #2
  \end{BoiteDfn}%
}

\newtcolorbox{BoitePrp}[1][]{
  enhanced, breakable,
  colback=bleuPropriete!6, colframe=bleuPropriete,
  boxrule=0pt, borderline west={3pt}{0pt}{bleuPropriete},
  arc=0pt, outer arc=0pt,
  left=14pt, right=14pt, top=10pt, bottom=10pt,
  before skip=12pt, after skip=12pt,
}
\newcommand{\Prp}[2][]{%
  \begin{BoitePrp}
    \refstepcounter{propriete}%
    \textbf{\color{bleuPropriete}Propriété~\arabic{propriete}}%
    \ifx\\#1\\\else\ --~\textbf{#1}\fi
    \par\smallskip #2
  \end{BoitePrp}%
}

\newtcolorbox{BoiteThm}[1][]{
  enhanced, breakable,
  colback=bordeauxTheoreme!6, colframe=bordeauxTheoreme,
  boxrule=0pt, borderline west={3pt}{0pt}{bordeauxTheoreme},
  arc=0pt, outer arc=0pt,
  left=14pt, right=14pt, top=10pt, bottom=10pt,
  before skip=12pt, after skip=12pt,
}
\newcommand{\Thm}[2][]{%
  \begin{BoiteThm}
    \refstepcounter{theoreme}%
    \textbf{\color{bordeauxTheoreme}Théorème~\arabic{theoreme}}%
    \ifx\\#1\\\else\ --~\textbf{#1}\fi
    \par\smallskip #2
  \end{BoiteThm}%
}

\newtcolorbox{BoiteExemple}[1][]{
  enhanced, breakable,
  colback=grisExemple!6, colframe=grisExemple,
  boxrule=0pt, borderline west={3pt}{0pt}{grisExemple},
  arc=0pt, outer arc=0pt,
  left=14pt, right=14pt, top=10pt, bottom=10pt,
  before skip=12pt, after skip=12pt,
}
\newcommand{\Exemple}[2][]{%
  \begin{BoiteExemple}
    \refstepcounter{exemple}%
    \textbf{\color{grisExemple}Exemple~\arabic{exemple}}%
    \ifx\\#1\\\else\ --~\textbf{#1}\fi
    \par\smallskip #2
  \end{BoiteExemple}%
}

\newtcolorbox{BoiteANoter}[1][]{
  enhanced, breakable,
  colback=ambreNoter!8, colframe=ambreNoter,
  boxrule=0pt, borderline west={3pt}{0pt}{ambreNoter},
  arc=0pt, outer arc=0pt,
  left=14pt, right=14pt, top=10pt, bottom=10pt,
  before skip=12pt, after skip=12pt,
}
\newcommand{\ANoter}[2][]{%
  \begin{BoiteANoter}
    \textbf{\color{ambreNoter}À noter}%
    \ifx\\#1\\\else\ --~\textbf{#1}\fi
    \par\smallskip #2
  \end{BoiteANoter}%
}

\newtcolorbox{BoiteAttention}[1][]{
  enhanced, breakable,
  colback=rougeAttention!7, colframe=rougeAttention,
  boxrule=0pt, borderline west={3pt}{0pt}{rougeAttention},
  arc=0pt, outer arc=0pt,
  left=14pt, right=14pt, top=10pt, bottom=10pt,
  before skip=12pt, after skip=12pt,
}
\newcommand{\Attention}[2][]{%
  \begin{BoiteAttention}
    \textbf{\color{rougeAttention}Attention}%
    \ifx\\#1\\\else\ --~\textbf{#1}\fi
    \par\smallskip #2
  \end{BoiteAttention}%
}

% ============================================================ %
% Cartes vidéos Yvan Monka : icône lecture (grande), titre cliquable,
% QR code cliquable (agrandi, 2.6cm de côté, pour rester scannable même
% imprimé en grand format). Une vidéo par ligne pleine largeur.
% ============================================================ %
\newcommand{\ClickableQR}[2]{\makebox[0pt][l]{\href{#2}{\makebox[#1][c]{\rule{0pt}{#1}}}}}
\newcommand{\VideoYvanLigne}[2]{%
  \noindent
  \makebox[0.1\linewidth]{$\vcenter{\hbox{\Large\faPlayCircle}}$}%
  \makebox[0.6\linewidth][l]{$\vcenter{\hbox{\parbox{0.58\linewidth}{\raggedright\large\href{#2}{#1}}}}$}%
  \makebox[0.3\linewidth]{$\vcenter{\hbox{\ClickableQR{2.2cm}{#2}\qrcode*[height=2.2cm, level=M]{#2}}}$}%
}
\newcommand{\VideosYvanMonka}[4]{{%
  \centering
  \textbf{Vidéos Yvan Monka}\par\smallskip
  \VideoYvanLigne{#1}{#2}

  \medskip

  \VideoYvanLigne{#3}{#4}
\par}}
\newcommand{\VideoYvanSolo}[2]{{%
  \centering
  \textbf{Vidéo Yvan Monka}\par\smallskip
  \VideoYvanLigne{#1}{#2}
\par}}

% Case à cocher (remplace \Square de bbding, non chargé ici).
\newcommand{\Case}{\raisebox{0pt}{\framebox(18,18){}}}

% Numéro cerclé simple (remplace \numeroCerne de CLT2627.sty).
\newcommand{\NumCercle}[1]{\textbf{(#1)}}

% Espace à compléter par l'élève : trait bas gris clair, largeur généreuse.
\newcommand{\Trou}[1][6cm]{\makebox[#1][l]{\textcolor{gray!55}{\rule[-3pt]{#1}{0.7pt}}}}

% Couleur des complétions (version complétée uniquement) : même vert un
% peu foncé que dans la version complétée classique, distinct du rouge
% \surligneReponse des mots clés du cours et du bleu des vecteurs donnés.
\definecolor{vertComplet}{RGB}{0,125,50}
\newcommand{\Complet}[1]{\textcolor{vertComplet}{#1}}
% Styles TikZ des constructions révélées : pointe Stealth agrandie pour le
% format grands caractères ; vecteurs intermédiaires (verts) en tirets,
% vecteur somme (rouge) en trait plein.
\tikzset{
  flecheComplet/.style={-{Stealth[length=12pt,width=9pt]},very thick},
  vecInter/.style={flecheComplet,vertComplet,dashed},
  vecSomme/.style={flecheComplet,red},
}
% Case cochée (réponse) : même gabarit que \Case, coche verte dedans.
\newcommand{\CaseOui}{\raisebox{0pt}{\framebox(18,18){\textcolor{vertComplet}{\small$\checkmark$}}}}

% ============================================================ %
% En-tête / pied de page
% ============================================================ %
\pagestyle{fancy}
\fancyhf{}
\renewcommand{\headrulewidth}{0pt}
\renewcommand{\footrulewidth}{0pt}
\fancyfoot[C]{\small Séquence 2 — Vecteurs — \thepage/\pageref{LastPage}}

\title{\vspace{-2cm}}
\date{}

\begin{document}

% ============================================================ %
% Titre — seul sur la page 1, saut de page avant la section 1 (comme
% dans la version classique, dont le titre n'est jamais mêlé au cours).
% ============================================================ %
\vspace*{3cm}
\SequenceTitre{0.7}{Géométrie - Séquence 2}{Vecteurs du plan}

\begin{center}
\textit{\Large\color{gray}Version complétée}
\end{center}

\clearpage

% ============================================================ %
% Section 1 — Définition d'un vecteur
% ============================================================ %
\section{Définition d'un vecteur}
\subsection{Translations et vecteurs}

\Dfn[Vecteur \texorpdfstring{$\Vj{AB}$}{AB}]{
Soient A et B deux points distincts du plan.\\
Le \textbf{vecteur} $\Vj{AB}$ est l'objet mathématique caractérisé par :
\begin{itemize}
\item sa \textbf{direction} : la droite (AB) ;
\item son \textbf{sens} : de A vers B ;
\item sa longueur, ou \textbf{norme}, notée $||\Vj{AB}||$ : la distance AB.
\end{itemize}
L'\textbf{origine} du vecteur $\Vj{AB}$ est A et son \textbf{extrémité} est B.

\begin{center}
\begin{tikzpicture}[x=1.6cm,y=1.6cm,>=Stealth,line width=1.4pt]
\draw[-latex,thick,blue] (0,0) node[below left]{A} -- (3,1.5) node[above right]{B};
\end{tikzpicture}
\end{center}
}

\Attention{On fera attention à ne pas confondre le vecteur $\Vj{AB}$ et le vecteur $\Vj{BA}$ (ils n'ont pas le même sens).}

\clearpage
\Dfn[Translation de vecteur \texorpdfstring{$\Vj{u}$}{u}]{
Soit $\Vj{u}$ un vecteur.\\
La \textbf{translation de vecteur $\Vj{u}$} est la transformation qui, à tout point A du plan, associe l'unique point A' tel que \textcolor{surligneReponse}{$\Vj{AA'} = \Vj{u}$}.

\begin{center}
\begin{tikzpicture}[yscale=0.9, xscale=1.3, line width=1.4pt]
    \node[left] at (1,0.7) {$\Vj{u}$};
\draw[->, line width=1pt] (-0.5,0)--(2.5,1) ;
\draw[-latex,thick,blue,dotted](1,-1) node[below left]{A}--(4,0) node[above right]{A'};
\end{tikzpicture}
\end{center}
}

En particulier, la translation de vecteur $\Vj{AB}$ transforme A en B (car $\Vj{AB}=\Vj{AB}$).

\clearpage
\Prp[Caractérisation par un parallélogramme]{
Soient A, B, C, D quatre points du plan.\\
Le point D est l'image de C par la translation de vecteur $\Vj{AB}$ \textbf{si et seulement si} \textcolor{surligneReponse}{ABDC est un parallélogramme}.

\begin{center}
\begin{tikzpicture}[yscale=0.9, xscale=1.3, line width=1.4pt]
\draw[-latex,thick,blue](0,0) node[below left]{A}--(3,1) node[above right]{B};
\draw[-latex,thick,blue](1,-1) node[below left]{C}--(4,0) node[above right]{D};
\draw (0,0)--(4,0);
\draw (1,-1)--(3,1);
\draw[red,thick] (1,0) node[scale=0.9]{|};
\draw[red,thick] (3,0) node[scale=0.9]{|};
\draw[red,thick] (1.5,-0.5) node[rotate=45,scale=0.9]{||};
\draw[red,thick] (2.5,0.5) node[rotate=45,scale=0.9]{||};
\end{tikzpicture}\\
{\small \textit{Le parallélogramme peut être aplati.}}
\end{center}

Cela équivaut aussi à dire que les segments [AD] et [BC] ont le \textcolor{surligneReponse}{même milieu}.
}

\clearpage
\Exemple{
Soit $ABCD$ un parallélogramme. Construire les points $E$, $F$, $G$ et $H$ tels que :
\begin{enumerate}
\item $E$ est l'image de $A$ par la translation de vecteur $\Vj{CB}$ ;
\item $F$ est l'image de $C$ par la translation de vecteur $\Vj{AB}$ ;
\item $G$ est l'image de $C$ par la translation de vecteur $\Vj{AD}$ ;
\item $H$ est l'image de $A$ par la translation de vecteur $\Vj{CD}$.
\end{enumerate}

\begin{center}
\begin{tikzpicture}[scale=1.1,x=1cm, y=1cm, >=Stealth, line width=1.3pt]
    \draw[thin, gray!50] (-3,-3) grid (6,5);
    \coordinate (A) at (0,0);
    \coordinate (B) at (2,1);
    \coordinate (C) at (3,3);
    \coordinate (D) at (1,2);
    \draw[thick] (A) -- (B) -- (C) -- (D) -- cycle;
    \fill (A) circle (3pt) node[below left] {$A$};
    \fill (B) circle (3pt) node[below right] {$B$};
    \fill (C) circle (3pt) node[above right] {$C$};
    \fill (D) circle (3pt) node[above left] {$D$};
    % Complétion : E = A + CB, F = C + AB, G = C + AD, H = A + CD
    % (CB = (-1;-2), AB = (2;1), AD = (1;2), CD = (-2;-1)).
    \coordinate (E) at (-1,-2);
    \coordinate (F) at (5,4);
    \coordinate (G) at (4,5);
    \coordinate (H) at (-2,-1);
    \draw[-latex,thick,vertComplet,dotted] (A) -- (E);
    \draw[-latex,thick,vertComplet,dotted] (C) -- (F);
    \draw[-latex,thick,vertComplet,dotted] (C) -- (G);
    \draw[-latex,thick,vertComplet,dotted] (A) -- (H);
    \fill[vertComplet] (E) circle (3pt) node[below right,vertComplet] {$E$};
    \fill[vertComplet] (F) circle (3pt) node[above right,vertComplet] {$F$};
    \fill[vertComplet] (G) circle (3pt) node[above,vertComplet] {$G$};
    \fill[vertComplet] (H) circle (3pt) node[below left,vertComplet] {$H$};
\end{tikzpicture}
\end{center}
}

\clearpage
\subsection{Égalité de vecteurs}

\Prp{Deux vecteurs sont \textbf{égaux} s'ils ont même direction, même sens et même norme.}

\ANoter{Un vecteur n'est pas caractérisé par son origine et son extrémité : il admet une \textbf{infinité de représentants}. Autrement dit, l'origine et l'extrémité d'un vecteur ne sont pas fixes.}

\Dfn{
Si $\Vj{AB}=\Vj{CD}=\Vj{EF}$, on dit que $\Vj{AB}$, $\Vj{CD}$ et $\Vj{EF}$
sont des \textbf{représentants} d'un même vecteur ($\Vj{u}$).

\begin{center}
\begin{tikzpicture}[scale=1.15, line width=1.3pt]
\draw[-latex,thick,blue] (0.4,-0.6)--(1.4,-1.6) node[midway, above right=-1pt]{$\Vj u$};
\draw[-latex,thick,blue] (0.2,1.6) node[above]{$A$}--(1.2,0.6) node[below right]{$B$};
\draw[-latex,thick,blue] (4.4,1.1) node[above]{$C$}--(5.4,0.1) node[below right]{$D$};
\draw[-latex,thick,blue] (2.6,-1.3) node[above]{$E$}--(3.6,-2.3) node[below right]{$F$};
\end{tikzpicture}
\end{center}
}

\clearpage
\Exemple{
Utiliser le quadrillage pour déterminer si les vecteurs suivants sont égaux : cocher, pour chaque paire, les critères vérifiés, puis conclure.

\begin{center}
\renewcommand{\arraystretch}{1.8}
\begin{tabular}{|l|c|c|c|c|}
\hline
 & \shortstack{même\\direction} & \shortstack{même\\sens} & \shortstack{même\\norme} & Conclusion \\
\hline
\NumCercle{1}\ $\Vj{AB}$ et $\Vj{EF}$ & \CaseOui & \CaseOui & \CaseOui & \makebox[3.5cm][l]{\Complet{Égaux}}\\
\hline
\NumCercle{2}\ $\Vj{DA}$ et $\Vj{DB}$ & \Case & \Case & \Case & \makebox[3.5cm][l]{\Complet{Non égaux}}\\
\hline
\NumCercle{3}\ $\Vj{ED}$ et $\Vj{BD}$ & \CaseOui & \Case & \CaseOui & \makebox[3.5cm][l]{\Complet{Non égaux}}\\
\hline
\NumCercle{4}\ $\Vj{AE}$ et $\Vj{BF}$ & \CaseOui & \CaseOui & \CaseOui & \makebox[3.5cm][l]{\Complet{Égaux}}\\
\hline
\end{tabular}
\end{center}

\begin{center}
\begin{tikzpicture}[x=0.8cm, y=0.8cm, line width=1.2pt]
    \draw[gray, line width=0.4pt,step=1] (0,0) grid (10,8);
    \path
        (2,7) node[circle, fill, inner sep=2.5pt, label=above: A] (A) {}
        (7,7) node[circle, fill, inner sep=2.5pt, label=above: B] (B) {}
        (9,4) node[circle, fill, inner sep=2.5pt, label=right: C] (C) {}
        (4,4) node[circle, fill, inner sep=2.5pt, label=left: D] (D) {}
        (1,1) node[circle, fill, inner sep=2.5pt, label=below:E] (E) {}
        (6,1) node[circle, fill, inner sep=2.5pt, label=below: F] (F) {};
\end{tikzpicture}
\end{center}
}

\clearpage
\subsection{Vecteurs particuliers}

\Dfn{
\begin{itemize}
\item la transformation qui laisse tous les points fixes est appelée translation de \textbf{vecteur nul} $\Vj{0}$.\\
La norme du vecteur nul est 0 mais sa direction et son sens ne sont pas définis.
\item on appelle \textbf{vecteur opposé} au vecteur $\Vj{AB}$ le vecteur noté $-\Vj{AB}$
associé à la translation qui transforme $B$ en $A$. \\
Autrement dit, $-\Vj{AB}=\Vj{BA}$.\\
Le vecteur $-\Vj{AB}$ a donc la même direction et la même norme que $\Vj{AB}$ mais est de sens opposé.
\end{itemize}
}

\Exemple{
\begin{center}
\begin{tikzpicture}[scale=1.5, line width=1.4pt]
\draw[-latex,thick] (0,0) node[left]{$A$}-- (2,1) node[right]{$B$} node[midway,rotate=27,above=1pt]{$\Vj{AB}$};
\draw[<-,thick] (2.7,0) -- (4.7,1)  node[midway,rotate=27,below=-1pt]{$-\Vj{AB}$};
\end{tikzpicture}
\end{center}
}

\clearpage
\subsection{Interprétation géométrique}

\Thm{
$ABCD$ est un parallélogramme \textbf{si et seulement si} \textcolor{surligneReponse}{$\Vj{AB}$ = $\Vj{DC}$}.

\medskip

Le parallélogramme peut être éventuellement « plat ».

\medskip

\textbf{Notation : } \textcolor{surligneReponse}{on note aussi que cette égalité entraîne $\Vj{AD}=\Vj{BC}$.}

\begin{center}
\begin{tikzpicture}[yscale=0.85, xscale=1.3, line width=2pt]
\draw[dashed] (1,1) -- (2,3);
\draw[dashed] (4,4) -- (3,2);
\draw[-{Stealth[length=6pt,width=5pt]}, line width=2pt] (2,3) -- (4,4);
\draw[-{Stealth[length=6pt,width=5pt]}, line width=2pt] (1,1) -- (3,2);
\node[anchor=east] at (1,1) {A};
\node[anchor=north west] at (3,2) {B};
\node[anchor=south east] at (2,3) {D};
\node[anchor=south west] at (4,4) {C};
\end{tikzpicture}
\end{center}
}

\clearpage
% ============================================================ %
% Section 2 — Somme de deux vecteurs
% ============================================================ %
\section{Somme de deux vecteurs}
\subsection{Définition et construction bout à bout}

\Dfn{
Soient $\Vj{u}$ et $\Vj{v}$ deux vecteurs. \\
La \textbf{somme} de $\Vj{u}$ et $\Vj{v}$ est le vecteur $\Vj{w}=\Vj{u}+\Vj{v}$ associé
à la translation résultant de l'enchaînement des translations de vecteurs $\Vj{u}$ et de vecteur $\Vj{v}$.
}

\clearpage
\textbf{\Large Méthode de construction bout à bout d'une somme de vecteurs}

\medskip

{\raggedcolumns
\begin{multicols}{3}
\begin{enumerate}[itemsep=3pt]
\item[] On veut construire la somme des vecteurs $\Vj{u}$ et $\Vj{v}$ :

\begin{tikzpicture}[scale=0.8]
\draw[thin, gray,opacity=0.7] (0,0) grid (6,6);
\draw[-{Stealth},thick,blue] (0,0) node[left]{}--(2,1) node[above]{} node[midway,rotate=30,above]{{\large $\Vj u$}};
\draw[-{Stealth},thick,blue] (2,5)--(4,2) node[right]{} node[midway,rotate=-60,above]{{\large $\Vj v$}};
\end{tikzpicture}

\columnbreak

\item Si l'origine du vecteur somme n'est pas imposée, on choisit un point M quelconque.
\item On construit N tel que $\Vj{MN}=\Vj{u}$.

\begin{tikzpicture}[scale=0.8]
\draw[thin, gray,opacity=0.7] (0,0) grid (6,6);
\fill[vertComplet] (1,4) circle (2.5pt);
\draw[vecInter] (1,4) node[left,vertComplet]{M}--(3,5) node[above,vertComplet]{N} node[midway,rotate=30,above]{{\large $\Vj u$}};
\end{tikzpicture}

\columnbreak

\item On construit P tel que $\Vj{NP}=\Vj{v}$. Alors $\Vj{u}+\Vj{v}=\Vj{MP}$.
\item On repère en \textcolor{red}{rouge} le vecteur somme et on n'oublie pas de le nommer.

\begin{tikzpicture}[scale=0.8]
\draw[thin, gray,opacity=0.7] (0,0) grid (6,6);
\fill[vertComplet] (1,4) circle (2.5pt);
\draw[vecInter] (1,4) node[left,vertComplet]{M}--(3,5) node[above,vertComplet]{N} node[midway,rotate=30,above]{{\large $\Vj u$}};
\draw[vecInter] (3,5)--(5,2) node[right,vertComplet]{P} node[midway,rotate=-60,above]{{\large $\Vj v$}};
% Vecteur somme en rouge, comme le demande l'étape 4 de la méthode.
\draw[vecSomme] (1,4)--(5,2) node[midway,rotate=-25,below=10pt]{{\large $\Vj u+\Vj v$}};
\end{tikzpicture}

\end{enumerate}
\end{multicols}
}

\ANoter{L'ordre dans une somme de vecteurs n'a pas d'importance :\\ $\Vj{u}+\Vj{v}=\Vj{v}+\Vj{u}$.}

\clearpage
\Exemple{
Construire le point $D$ tel que $\Vj{ED}=\Vj{BA}+\Vj{CA}$.

\begin{center}
\begin{tikzpicture}[scale=0.85]
\draw[thin, gray,opacity=0.7] (-2,0) grid (5,6);
\draw [ ] (0,2) node[above left]{A} -- ++(-4pt,-4pt) -- ++(8pt,8pt) ++(-8pt,0) -- ++(8pt,-8pt);
\draw [ ] (3,3) node[above left]{B} -- ++(-4pt,-4pt) -- ++(8pt,8pt) ++(-8pt,0) -- ++(8pt,-8pt);
\draw [ ] (1,1) node[above left]{C} -- ++(-4pt,-4pt) -- ++(8pt,8pt) ++(-8pt,0) -- ++(8pt,-8pt);
\draw [ ] (3,5) node[below right]{E} -- ++(-4pt,-4pt) -- ++(8pt,8pt) ++(-8pt,0) -- ++(8pt,-8pt);
% Complétion : E + BA = (0;4), puis + CA = (-1;5) = D ; vecteur somme
% ED en rouge, comme dans la méthode.
\draw[vecInter] (3,5)--(0,4) ;
\draw[vecInter] (0,4)--(-1,5) ;
\draw[vecSomme] (3,5)--(-1,5) node[midway,above]{$\Vj{ED}$};
\draw [vertComplet] (-1,5) node[above right]{D} -- ++(-4pt,-4pt) -- ++(8pt,8pt) ++(-8pt,0) -- ++(8pt,-8pt);
\end{tikzpicture}
\end{center}
}

\clearpage
\begin{BoiteExemple}
Construire le point $F$ tel que $\Vj{GF}=\Vj{AB}-\Vj{AC}$.

\begin{center}
\begin{tikzpicture}[scale=0.85]
\draw[thin, gray,opacity=0.7] (-1,0) grid (5,5);
\draw [ ] (0,2) node[above left]{A} -- ++(-4pt,-4pt) -- ++(8pt,8pt) ++(-8pt,0) -- ++(8pt,-8pt);
\draw [ ] (3,3) node[above right]{B} -- ++(-4pt,-4pt) -- ++(8pt,8pt) ++(-8pt,0) -- ++(8pt,-8pt);
\draw [ ] (1,1) node[above left]{C} -- ++(-4pt,-4pt) -- ++(8pt,8pt) ++(-8pt,0) -- ++(8pt,-8pt);
\draw [ ] (1,0) node[above left]{G} -- ++(-4pt,-4pt) -- ++(8pt,8pt) ++(-8pt,0) -- ++(8pt,-8pt);
% Complétion : G + AB = (4;1), puis - AC = + CA = (3;2) = F.
\draw[vecInter] (1,0)--(4,1) ;
\draw[vecInter] (4,1)--(3,2) ;
\draw[vecSomme] (1,0)--(3,2) node[midway,below right=-1pt]{$\Vj{GF}$};
\draw [vertComplet] (3,2) node[above right]{F} -- ++(-4pt,-4pt) -- ++(8pt,8pt) ++(-8pt,0) -- ++(8pt,-8pt);
\end{tikzpicture}
\end{center}
\end{BoiteExemple}

\clearpage
\subsection{Relation de Chasles}

\Prp[Relation de Chasles]{
Soient $A$ et $B$ deux points du plan. Pour tout point $M$, on a :
\[
\Vj{AB}=\Vj{AM}+\Vj{MB}.
\]
}

\ANoter{Aller d'Annoville à Brainville, puis de Brainville à Coutances, c'est finalement aller d'Annoville à Coutances.}

\clearpage
\Exemple{
Simplifier les expressions en utilisant la relation de Chasles :
\begin{enumerate}
\item $\Vj{TR}+\Vj{RP}$\\[4pt]
$= \Complet{\Vj{TP}}$
\item $\Vj{GT}+\Vj{AG}$\\[4pt]
$= \Complet{\Vj{AG}+\Vj{GT} = \Vj{AT}}$
\item $\Vj{AB}-\Vj{AC}-\Vj{CB}$\\[4pt]
$= \Complet{\Vj{AB}+\Vj{CA}+\Vj{BC} = \Vj{AB}+\Vj{BC}+\Vj{CA} = \Vj{AC}+\Vj{CA} = \Vj{AA} = \Vj{0}}$
\item $\Vj{BC}-\Vj{BA}+\Vj{BD}-\Vj{BC}$\\[4pt]
$= \Complet{\Vj{BC}+\Vj{AB}+\Vj{BD}+\Vj{CB} = \Vj{AB}+\Vj{BD}+\Vj{BC}+\Vj{CB} = \Vj{AD}+\Vj{BB} = \Vj{AD}}$
\end{enumerate}
}

\clearpage
\subsection{Méthode du parallélogramme}

\textbf{\Large Méthode de construction par parallélogramme d'une somme de vecteurs}

\medskip

{\raggedcolumns
\begin{multicols}{3}
\begin{enumerate}[itemsep=3pt]
\item[] On veut construire la somme des vecteurs $\Vj{u}$ et $\Vj{v}$ ayant comme origine commune le point M.

\begin{tikzpicture}[scale=0.8]
\draw[thin, gray,opacity=0.7] (0,0) grid (6,6);
\draw[-{Stealth},thick,blue] (1,1)--(3,2) node[midway,rotate=27,below]{{\large $\Vj u$}};
\draw[-{Stealth},thick,blue] (1,1)--(2,4) node[midway,rotate=72,above]{{\large $\Vj v$}};
\fill (1,1) circle (2.5pt);
\node[anchor=north east] at (1,1) {M};
\end{tikzpicture}

\columnbreak

\item On repère N et P tels que $\Vj{MN}=\Vj{u}$ et $\Vj{MP}=\Vj{v}$.

\begin{tikzpicture}[scale=0.8]
\draw[thin, gray,opacity=0.7] (0,0) grid (6,6);
\draw[-{Stealth},thick,blue] (1,1)--(3,2) node[midway,rotate=27,below]{{\large $\Vj u$}};
\draw[-{Stealth},thick,blue] (1,1)--(2,4) node[midway,rotate=72,above]{{\large $\Vj v$}};
\fill (1,1) circle (2.5pt);
\node[anchor=north east] at (1,1) {M};
\fill[vertComplet] (3,2) circle (2.5pt);
\node[anchor=north west,vertComplet] at (3,2) {N};
\fill[vertComplet] (2,4) circle (2.5pt);
\node[anchor=south east,vertComplet] at (2,4) {P};
\end{tikzpicture}

\columnbreak

\item On construit Q tel que MNQP soit un parallélogramme. Alors $\Vj{u}+\Vj{v}=\Vj{MQ}$.

\begin{tikzpicture}[scale=0.8]
\draw[thin, gray,opacity=0.7] (0,0) grid (6,6);
\draw[-{Stealth},thick,blue] (1,1)--(3,2) node[midway,rotate=27,below]{{\large $\Vj u$}};
\draw[-{Stealth},thick,blue] (1,1)--(2,4) node[midway,rotate=72,above]{{\large $\Vj v$}};
\fill (1,1) circle (2.5pt);
\node[anchor=north east] at (1,1) {M};
\fill[vertComplet] (3,2) circle (2.5pt);
\node[anchor=north west,vertComplet] at (3,2) {N};
\fill[vertComplet] (2,4) circle (2.5pt);
\node[anchor=south east,vertComplet] at (2,4) {P};
\draw[dotted,thick,vertComplet] (3,2)--(4,5);
\draw[dotted,thick,vertComplet] (2,4)--(4,5);
\fill[vertComplet] (4,5) circle (2.5pt);
\node[anchor=south west,vertComplet] at (4,5) {Q};
% Étiquette posée à droite de Q, hors du parallélogramme (le long de [MQ]
% elle chevauchait soit v, soit le côté pointillé [NQ]).
\draw[vecSomme] (1,1)--(4,5);
\node[red,anchor=west] at (4.15,4.1) {{\large $\Vj u+\Vj v$}};
\end{tikzpicture}

\end{enumerate}
\end{multicols}
}

\clearpage
% ============================================================ %
% Section 3 — Multiplication par un scalaire
% ============================================================ %
\section{Multiplication par un scalaire}

\Exemple[introductif]{
\begin{enumerate}
	\item $\Vj{u}$ est le vecteur
			\begin{itemize}
				\item de même direction que $\Vj{v}$ ;
				\item de même sens que $\Vj{v}$ ;
				\item de norme $3$ fois celle de $\Vj{v}$.
			\end{itemize}
			On écrit alors : $\Complet{\Vj{u} = 3\Vj{v}}$
	\item $\Vj{j}$ est le vecteur
			\begin{itemize}
				\item de même direction que $\Vj{i}$ ;
				\item de sens opposé à $\Vj{i}$ ;
				\item de norme deux fois celle de $\Vj{i}$.
			\end{itemize}
			On écrit alors : $\Complet{\Vj{j} = -2\Vj{i}}$
\end{enumerate}

\begin{center}
\begin{tikzpicture}[scale=0.75]
\draw[thin,gray,opacity=0.7] (0,0) grid  (12,12);
\draw[->,ultra thick] (3,5)--(1,7) node[midway,rotate=-45,above]{$\Vj v$};
\draw[->,ultra thick,vertComplet] (7,5)--(1,11) node[midway,rotate=-45,above]{$\Vj u$};
\draw[->,ultra thick] (7,9)--(9,11) node[midway,rotate=45,above]{$\Vj i$};
\draw[->,ultra thick,vertComplet] (11,5)--(7,1) node[midway,rotate=45,above]{$\Vj j$};
\end{tikzpicture}
\end{center}
}

\clearpage
\Dfn[Produit d'un vecteur par un nombre réel]{
Soit $\Vj{u}$ un vecteur non nul et soit $k$ un réel non nul.\\
Alors le vecteur {\Large $\Vj{v}= \textcolor{red}{k}\Vj{u}$} est :

\begin{center}
\renewcommand{\arraystretch}{1.8}
\begin{tabular}{|p{10.5cm}|p{10.5cm}|}
\hline
\centering si $\textcolor{red}{k}>0$ & \centering si $\textcolor{red}{k}<0$\tabularnewline
\hline
de même direction que $\Vj{u}$  & de même direction que $\Vj{u}$ \\
de même sens que $\Vj{u}$ & de sens inverse de $\Vj{u}$\\
de norme $k$ fois celle de $\Vj{u}$ & de norme $(-k)$ fois celle de $\Vj{u}$\\
\hline
\end{tabular}
\end{center}
}

\clearpage
\Exemple{
\begin{enumerate}
\item Exprimer le vecteur $\Vj{AC}$ en fonction du vecteur $\Vj{AB}$ dans chacun des cas suivants :

\begin{center}
\begin{tabular}{p{0.46\linewidth} p{0.46\linewidth}}
\begin{tikzpicture}[x=2cm, y=1.3cm]
    \draw (-0.5,0) -- (2.5,0);
    \foreach \x in {-0.5,0,...,2.5} {
        \draw (\x, 0.08) -- (\x, -0.08);
    }
    \node[ptcross] at (0,0) {};
    \node[anchor=south,yshift=2pt] at (0,0) {$A$};
    \node[ptcross] at (1,0) {};
    \node[anchor=south,yshift=2pt] at (1,0) {$B$};
    \node[ptcross] at (2,0) {};
    \node[anchor=south,yshift=2pt] at (2,0) {$C$};
\end{tikzpicture}
&
\begin{tikzpicture}[x=2cm, y=1.3cm]
    \draw (-2,0) -- (2,0);
    \foreach \x in {-2,-1.5,...,2} {
        \draw (\x, 0.08) -- (\x, -0.08);
    }
    \node[ptcross] at (1,0) {};
    \node[anchor=south,yshift=2pt] at (1,0) {$A$};
    \node[ptcross] at (1.5,0) {};
    \node[anchor=south,yshift=2pt] at (1.5,0) {$B$};
    \node[ptcross] at (-1.5,0) {};
    \node[anchor=south,yshift=2pt] at (-1.5,0) {$C$};
\end{tikzpicture}\\[6pt]
$\Vj{AC} = \Complet{2\Vj{AB}}$ & $\Vj{AC} = \Complet{-5\Vj{AB}}$
\end{tabular}
\end{center}

% Rédaction du raisonnement (reprise de la version complétée classique :
% direction, sens, longueurs en graduations, rapport, conclusion), en
% deux colonnes sous les figures. En \small et sans minipage imbriquée
% (qui rétrécirait encore les colonnes) : la page doit rester la même que
% côté élève, une page par Exemple.
\begin{center}
\begin{tabular}{@{}p{0.46\linewidth}@{\hspace{0.04\linewidth}}p{0.46\linewidth}@{}}
\RaggedRight\Complet{Même direction, même sens ; $AC=4$ et $AB=2$ graduations, donc $AC=2\times AB$, soit $\Vj{AC}=2\,\Vj{AB}$.}
&
\RaggedRight\Complet{Même direction, sens opposé ; $AC=5$ et $AB=1$ graduation, donc $AC=5\times AB$, soit $\Vj{AC}=-5\,\Vj{AB}$.}
\end{tabular}
\end{center}
\end{enumerate}
}

\clearpage
\begin{BoiteExemple}
\begin{enumerate}
\setcounter{enumi}{1}
\item Dans chacun des cas suivants, représenter le vecteur $\Vj{AC}$ tel que :

\begin{center}
\begin{tabular}{p{0.46\linewidth} p{0.46\linewidth}}
$\Vj{AC} = -3\Vj{AB}$ & $\Vj{AC} = \dfrac{1}{3}\Vj{AB}$ \\[8pt]
\begin{tikzpicture}[x=1.6cm, y=1.3cm]
    \draw (-2,0) -- (2.5,0);
    \foreach \x in {-2,-1.5,...,2.5} {
        \draw (\x, 0.08) -- (\x, -0.08);
    }
    \node[ptcross] at (0,0) {};
    \node[anchor=south,yshift=2pt] at (0,0) {$A$};
    \node[ptcross] at (0.5,0) {};
    \node[anchor=south,yshift=2pt] at (0.5,0) {$B$};
    % Complétion : AB = 1 graduation, donc C à 3 graduations de A, à
    % l'opposé de B.
    \node[ptcross,vertComplet] at (-1.5,0) {};
    \node[anchor=south,yshift=2pt,vertComplet] at (-1.5,0) {$C$};
    \draw[flecheComplet,vertComplet] (0,0) -- (-1.5,0) node[midway,below=4pt,font=\small]{$\Vj{AC}$};
\end{tikzpicture}
&
\begin{tikzpicture}[x=1.6cm, y=1.3cm]
    \draw (-1.5,0) -- (2.5,0);
    \foreach \x in {-1.5,-1,...,2.5} {
        \draw (\x, 0.08) -- (\x, -0.08);
    }
    \node[ptcross] at (0,0) {};
    \node[anchor=south,yshift=2pt] at (0,0) {$A$};
    \node[ptcross] at (1.5,0) {};
    \node[anchor=south,yshift=2pt] at (1.5,0) {$B$};
    % Complétion : AB = 3 graduations, donc C à 1 graduation de A, du
    % côté de B.
    \node[ptcross,vertComplet] at (0.5,0) {};
    \node[anchor=south,yshift=2pt,vertComplet] at (0.5,0) {$C$};
    \draw[flecheComplet,vertComplet] (0,0) -- (0.5,0) node[midway,below=4pt,font=\small]{$\Vj{AC}$};
\end{tikzpicture}
\end{tabular}
\end{center}

% Rédaction de la construction (même canevas que la question 1 : sens
% d'après le signe de k, longueur d'après |k|), en \small pour tenir sur
% la même page que côté élève.
\begin{center}
\begin{tabular}{@{}p{0.46\linewidth}@{\hspace{0.04\linewidth}}p{0.46\linewidth}@{}}
\RaggedRight\Complet{Comme $-3<0$, sens opposé à $\Vj{AB}$ ; $AB=1$~graduation, donc $AC=3$~graduations.}
&
\RaggedRight\Complet{Comme $\dfrac13>0$, même sens que $\Vj{AB}$ ; $AB=3$~graduations, donc $AC=1$~graduation.}
\end{tabular}
\end{center}
\end{enumerate}
\end{BoiteExemple}

\clearpage
\Exemple{
\begin{enumerate}
\item Placer les points H, I et J tels que :\\
(a) $\Vj{AH}=\dfrac53\Vj{AB}$ \quad (b) $\Vj{BI}=\dfrac23\Vj{BC}$ \quad (c) $\Vj{DJ}=\dfrac58\Vj{DC}$

\begin{center}
\begin{tikzpicture}[scale=1]
\draw[thin,gray,opacity=0.5] (0,0) grid (10,6);
\draw[very thick,magenta] (3,5) -- (6,5) -- (9,2) -- (1,1) -- cycle;
\node[ptcross] at (3,5) {}; \node[anchor=south] at (3,5) {A};
\node[ptcross] at (4,5) {}; \node[anchor=south] at (4,5) {E};
\node[ptcross] at (6,5) {}; \node[anchor=south] at (6,5) {B};
\node[ptcross] at (9,2) {}; \node[anchor=west] at (9,2) {C};
\node[ptcross] at (1,1) {}; \node[anchor=north east] at (1,1) {D};
\node[ptcross] at (2,3) {}; \node[anchor=east] at (2,3) {F};
% H = A + 5/3 AB = (8 ; 5) ; I = B + 2/3 BC ; J = D + 5/8 DC = (1+5 ; 1+5/8).
\node[ptcross,vertComplet] at (8,5) {}; \node[anchor=south,vertComplet] at (8,5) {H};
\node[ptcross,vertComplet] at (8,3) {}; \node[anchor=west,vertComplet] at (8,3) {I};
\node[ptcross,vertComplet] at (6,1.625) {}; \node[anchor=south,vertComplet] at (6,1.625) {J};
\end{tikzpicture}
\end{center}
\end{enumerate}
}

\clearpage
\begin{BoiteExemple}
\begin{enumerate}
\setcounter{enumi}{1}
\item Construire dans la figure ci-dessous des représentants des vecteurs $2\Vj{u}$, $-\Vj{u}$ et $\dfrac12\Vj{u}$ :

\begin{center}
\begin{tikzpicture}[scale=1.15, line width=1.3pt]
\draw[thin,gray,opacity=0.7] (-3,-1) grid (6,3);
\draw[->,thick,blue] (0,0) -- (3,1) node[midway,rotate=20,above]{$\Vj{u}$};
\draw[->,thick,vertComplet] (-2,-1) -- (4,1) node[midway,rotate=20,below]{$2\Vj{u}$};
\draw[->,thick,vertComplet] (2,3) -- (-1,2) node[midway,rotate=20,above]{$-\Vj{u}$};
\draw[->,thick,vertComplet] (-3,2) -- (-1.5,2.5) node[midway,rotate=20,above]{$\frac12\Vj{u}$};
\end{tikzpicture}
\end{center}
\end{enumerate}
\end{BoiteExemple}

\clearpage
% ============================================================ %
% Section 4 — Vecteurs colinéaires
% ============================================================ %
\section{Vecteurs colinéaires}

\Dfn[Vecteurs colinéaires]{
Deux vecteurs $\Vj{u}$ et $\Vj{v}$ sont dits \textbf{colinéaires} s'il existe un réel $k$ tel que $\Vj{v}=k\Vj{u}$ ou $\Vj{u}=k\Vj{v}$.
}

\ANoter{
\begin{itemize}
\item deux vecteurs sont colinéaires si et seulement s'ils ont \textbf{même direction}.
\item le vecteur nul est colinéaire à n'importe quel vecteur. En effet, $\Vj{0}=0\Vj{u}$.
\end{itemize}
}

\clearpage
\Exemple{
\begin{center}
\begin{tikzpicture}[scale=1.3, line width=1.3pt]
\draw[thin,gray,opacity=0.5] (0,0) grid (6,4);
\draw[->,thick,blue] (0.5,0.5) -- (1.5,1) node[midway,rotate=27,above]{$\Vj u$};
\draw[->,thick,red] (5,3) -- (3,2) node[midway,rotate=27,above]{$\Vj v$};
\end{tikzpicture}
\end{center}
$\Vj{v} = \Complet{-2}\,\Vj{u}$

\Complet{Donc $\Vj{u}$ et $\Vj{v}$ sont colinéaires.}
}

\ANoter[À retenir]{La notion de colinéarité est fortement liée à celle de parallélisme et d'alignement.}

\clearpage
\Prp[Colinéarité et parallélisme]{
Soient A, B, C et D quatre points.

\medskip

Les vecteurs $\Vj{AB}$ et $\Vj{CD}$ sont colinéaires si et seulement si les droites (AB) et (CD) sont \textcolor{surligneReponse}{parallèles (ou confondues)}.

\begin{center}
\begin{tikzpicture}[scale=1.1, line width=1.6pt]
\draw[thin,gray,opacity=0.5] (0,0) grid (8,5);
\draw[|-{Stealth[length=10pt,inset=2pt]},very thick,blue] (1,2) node[below=2pt]{A} --(2,4) node[above=2pt]{B} node[midway,rotate=60,above=2pt]{$\Vj{AB}$};
\draw[dotted,red, thick](0,0) -- (1,2);
\draw[dotted,red, thick](2,4) --(2.5,5) ;
\draw[|-{Stealth[length=10pt,inset=2pt]},very thick,blue] (6,5) node[below right]{C} --(4,1) node[below=2pt]{D} node[midway,rotate=60,above=2pt]{$\Vj{CD}$};
\draw[dotted,red](3.5,0) -- (4,1);
\end{tikzpicture}
\end{center}

\textbf{Rédaction : }

\medskip

\Complet{Ici, on a : $\Vj{CD} = -2 \  \Vj{AB}$ donc $\Vj{AB}$ et $\Vj{CD}$ sont colinéaires.}

\medskip

\Complet{Ainsi : les droites $(AB)$ et $(CD)$ sont parallèles.}
}

\clearpage
\Prp[Colinéarité et alignement]{
Soient A, B et C trois points distincts.

\medskip

Les points A, B et C sont alignés si et seulement si les vecteurs $\Vj{AB}$ et $\Vj{AC}$ sont \textcolor{surligneReponse}{colinéaires}.

\begin{center}
\begin{tikzpicture}[scale=1.1, line width=1.6pt]
\draw[thin,gray,opacity=0.5] (0,1) grid (7,4);
\draw[opacity=0.8](0,1.8) --(7,3.2) ;
\draw[|->,very thick,red] (1,2) node[below=2pt]{A} --(6,3) node[below=2pt]{C} node[midway,rotate=10,below]{$\Vj{AC}$};
\draw[->, thick,blue] (1,2) --(4,2.6) node[above=2pt]{B} node[midway,rotate=10,above]{$\Vj{AB}$};
\end{tikzpicture}
\end{center}

\textbf{Rédaction : }

\medskip

\Complet{Ici, on a : $\Vj{AB} = \dfrac35 \  \Vj{AC}$ donc $\Vj{AB}$ et $\Vj{AC}$ sont colinéaires.}

\medskip

\Complet{Ainsi : les points A, B et C sont alignés.}
}

\end{document}
